\documentclass[10pt]{article}
\usepackage[T1]{fontenc}
\usepackage[utf8]{inputenc}
\usepackage[french]{babel}
\usepackage{amsmath,amssymb,mathrsfs,mathpazo}
\usepackage{geometry}
\geometry{a4paper,margin=2cm}
\usepackage{graphicx,array,tabularx,tabularray,colortbl,multirow,multicol,booktabs,enumitem,pifont,ulem,fancybox,diagbox,cancel}
\usepackage{pstricks,pstricks-add,pst-plot,pst-tree,pst-node,pst-text,pst-func,pst-math,pst-bspline,pst-3dplot,pst-eucl,pst-solides3d}
\let\pscalculate\undefined
\usepackage{pst-bezier}
\usepackage{tikz,pgfplots,tkz-tab}
\usetikzlibrary{arrows,arrows.meta,patterns,calc,positioning,trees,decorations.markings,shapes.geometric,shapes.misc,shapes.arrows,intersections,decorations.pathmorphing,decorations.pathreplacing,backgrounds,shadows,babel}
\pgfplotsset{compat=1.18}
\usepgfplotslibrary{fillbetween}
\usepackage[np]{numprint}
\usepackage{listings,makecell,pas-tableur,twoopt}
\let\barre\undefined
\usepackage[pstricks]{tablvar}
\definecolor{darkgreen}{RGB}{0,100,0}
\providecommand{\red}{\color{red}}
\providecommand{\blue}{\color{blue}}
\providecommand{\green}{\color{green}}
\providecommand{\black}{\color{black}}
\providecommand{\white}{\color{white}}
\providecommand{\gray}{\color{gray}}
\providecommand{\cyan}{\color{cyan}}
\providecommand{\magenta}{\color{magenta}}
\providecommand{\yellow}{\color{yellow}}
\providecommand{\violet}{\color{violet}}
\providecommand{\orange}{\color{orange}}
\providecommand{\euro}{\texteuro}
\providecommand{\EUR}{\texteuro}
\providecommand{\ieme}{\textsuperscript{e}}
\providecommand{\degres}{\ensuremath{^\circ}}
\providecommand{\degree}{\ensuremath{^\circ}}
\providecommand{\vv}[1]{\overrightarrow{#1}}
\providecommand{\Coord}[1]{#1}
\providecommand{\decosix}{}
\providecommand{\starredbullet}{\ensuremath{\bullet}}
\pagestyle{empty}
\setlength{\parindent}{0pt}
\usepackage{adjustbox}\DeclareUnicodeCharacter{2236}{:}\DeclareUnicodeCharacter{2190}{\ensuremath{\leftarrow}}
\providecommand{\R}{\mathbb{R}}
\providecommand{\N}{\mathbb{N}}
\providecommand{\Z}{\mathbb{Z}}
\providecommand{\Q}{\mathbb{Q}}
\providecommand{\D}{\mathbb{D}}
\providecommand{\C}{\mathbb{C}}
\providecommand{\e}{{\mathrm{e}}}
\providecommand{\vect}[1]{\overrightarrow{#1}}
\providecommand{\vectt}[1]{\overrightarrow{\mathrm{#1}}}
\providecommand{\vvt}[1]{\overrightarrow{\mathrm{#1}}}
\providecommand{\np}[2][]{{#2~#1}}
\providecommand{\vv}[1]{{\overrightarrow{#1}}}
\renewcommand{\d}{{\mathrm{d}}}
\providecommand{\degres}{{^\circ}}
\providecommand{\degre}{{^\circ}}
\providecommand{\blue}{\color{blue}}
\providecommand{\red}{\color{red}}
\providecommand{\enskip}{\quad}
\providecommand{\llbracket}{\lbrack\!\lbrack}
\providecommand{\rrbracket}{\rbrack\!\rbrack}
\providecommand{\euro}{{\text{€}}}
\providecommand{\no}{{\text{n°}}}
\providecommand{\interval}[2]{{[#1;#2]}}
\providecommand{\intervFO}[2]{{[#1;#2[}}
\providecommand{\dotfill}{\ldots}

\providecommand{\ds}{\displaystyle}
\providecommand{\textding}[1]{#1}
\providecommand{\Oij}{\(\mathrm{O};\vec{\imath},\vec{\jmath}\)}
\providecommand{\Oijk}{\(\mathrm{O};\vec{\imath},\vec{\jmath},\vec{k}\)}
\providecommand{\Ouv}{\(\mathrm{O};\vec{u},\vec{v}\)}
\providecommand{\intervFF}[2]{\ensuremath{[#1;#2]}}
\renewcommand{\euro}{\ifmmode\text{\texteuro}\else\texteuro\fi}
\DeclareUnicodeCharacter{00B0}{\ifmmode^{\circ}\else\textdegree\fi}
\DeclareUnicodeCharacter{20AC}{\ifmmode\text{\texteuro}\else\texteuro\fi}
% Énoncés : Éducation nationale. Transcription et corrigés : APMEP, Denis Vergès et contributeurs. Découpage et classement : CoopMaths, Sébastien Lozano, Éric Elter et contributeurs. Adaptation de la présentation : Maths Explorer.
% CC BY-SA : https://coopmaths.fr/www/about/
% https://coopmaths.fr/alea/static/bac/2021/tex/bac_2021_06_etrangers_6_cor.tex
\begin{document}
\noindent\textit{APMEP / CoopMaths. CC BY-SA. Présentation adaptée par Maths Explorer.}\par\medskip
\textbf{Commun à tous les candidats}

\medskip














\textbf{Question 1 :}

On considère la fonction $g$ définie sur $]0~;~+\infty[$  par $g(x)= x^2+ 2x - \dfrac{3}{x}$.







Si $T$ est la tangente à la courbe représentative de $g$ au point d'abscisse 1, on sait que :

$M(x~;~y) \in T \iff y - g(1) = g'(1)(x - 1)$.

\starredbullet $g$ est dérivable sur $]0~;~+\infty[$ et sur cet intervalle : $g'(x) = 2x + 2 + \dfrac{3}{x^2}$

On a donc $g'(1) = 2 + 2 + \dfrac{3}{1} = 7$ ;

\starredbullet D'autre part $g(1) = 1^2 + 2 \times 1 - \dfrac{3}{1} = 0$.

On a donc : $M(x~;~y) \in T \iff y - 0 = 7(x - 1) \iff y = 7x - 7$. Réponse \textbf{a.}

\medskip

\textbf{Question 2 :}









Pour $n$ assez grand, on a $n \ne 0$, donc $v_n = \dfrac{3}{1 + \frac{2}{n}}$

Comme $\displaystyle\lim_{n \to + \infty} \dfrac{3}{n} = 0$, on a $\displaystyle\lim_{n \to + \infty} = 3$. Réponse \textbf{b.}

\medskip

\textbf{Question 3 :}








Si $X$ est la variable aléatoire correspondant au nombre de boules noires tirées, l'expérience correspond à une loi de Bernoulli  et $X$ suit une loi binomiale de paramètres $n = 10$ et $p = 0,6$.

On sait que $p(X = 4) = \binom{10}{4} \times 0,6^4 \times \left(1 - 0,4^6\right) \approx \np{0,11147} \approx \np{0,1115}$ à $10^{-4}$ près. Réponse \textbf{c.~~}
\medskip

\textbf{Question 4 :}









On a pour $x \ne 0$, \, $f(x) = x\left(3\dfrac{\e^x}{x}  - 1\right)$.

Or on sait que $\displaystyle\lim_{x \to + \infty} \dfrac{\text{e}^x}{x} = + \infty$, donc 
$\displaystyle\lim_{x \to + \infty} \dfrac{\text{e}^x}{x} - 1 = + \infty$ et par produit de limites
 
$\displaystyle\lim_{x \to + \infty} f(x) = + \infty$ Réponse \textbf{b.}
\medskip

\textbf{Question 5 :}














Le nombre de codes différents est donc $36^8 = \np{2821109907456}$.

Il faut donc $\dfrac{\np{2821109907456}}{10^8} \approx \np{28211}$~(s) ou $\approx 478,2$~(min) ou $\approx 7,8$~(h). Réponse \textbf{b.}

\bigskip

\end{document}
